\chapter*{后记}
\addcontentsline{toc}{chapter}{后记}
\markboth{后记}{后记}

写到微分方程时，第一章的数列又有了新的用途。递推由变化规律产生，有限步骤给出近似，误差估计说明这些近似怎样接近一个函数。极限、导数和积分在同一个问题中发挥作用，前面的知识也因此获得了进一步的意义。

这样的联系需要充分的准备。求导与积分起初分别研究局部的变化率与区间上的累积量；经过极限与连续性的论证，二者才在微积分基本定理中联系起来。借助原函数计算累积量，后来成为熟练的运算；回望它的建立过程，仍能体会这一联系的分量。

这些理论早已成熟，需要重新完成的，是把它们建立在读者已有的经验上。编写中，许多调整落在概念的次序、必要的准备和推理的接续上。我在意这些安排，因为读者正是沿着具体的一步步向前。眼前的推导有来处，学过的内容在后来仍有用途，零散的经验才逐渐成为可以依靠的知识。

我相信，高中生能够沿着这样的道路走向更深的理解。讲义提供必要的准备与清楚的依据，也把能够独立完成的部分留给读者。自己补出的一步推理、自己发现的一处联系，使知识更为可靠。再遇到陌生问题，便多了一些可以从容出发的地方。

筹备概率统计一章时，前面的工具已经能够支持许多顺畅的延伸。此时更需要判断，哪些关系值得展开，一个问题研究到哪里便可以结束。各章都需要这样的控制，工具越充分，选择越显重要。作者知道更多，可以用来更准确地安排读者此刻需要理解的内容，讲义也由此形成自己的取向。

读者也有自己的选择。有些内容此时就能用上，有些认识要在以后的学习中获得。再次翻到原来的一页，文字或许没有改变，能够看出的关系却多了一些。我希望这份讲义经得起这样的重读。曾经需要反复解释的结论，也能由读者自己清楚地讲出来。

当一个结论能够由自己推得，一个陌生问题能够从已有概念出发理解，数学便逐渐成为一种可以依靠的思考方式。我把自己的判断落实在文字里，也期待读者形成自己的判断，并由此继续前行。《从数字到数学》，如是而已。

\begin{flushright}
2026年11月14日
\end{flushright}
